%% =====================================================================
%%  macros.tex  --  theorem environments and maths shortcuts
%% =====================================================================

%% ---- Theorem styles --------------------------------------------------
%% (Palatino has no bold small caps, so a head is bold OR small caps.)
\newtheoremstyle{statement}   % claims you cite later: bold head, italic body
  {9pt plus 2pt minus 2pt}{5pt plus 2pt minus 2pt}{\itshape}{0pt}{\bfseries}{.}{0.5em}{}
\newtheoremstyle{explain}     % definitions, examples: bold head, upright body
  {9pt plus 2pt minus 2pt}{5pt plus 2pt minus 2pt}{\normalfont}{0pt}{\bfseries}{.}{0.5em}{}
\newtheoremstyle{aside}       % problems, remarks: small-caps head, upright body
  {9pt plus 2pt minus 2pt}{5pt plus 2pt minus 2pt}{\normalfont}{0pt}{\scshape}{.}{0.5em}{}

%% One shared counter, so "Theorem 1.3" and "Lemma 1.3" can't both exist.
\theoremstyle{statement}
\newtheorem{theorem}{Theorem}[section]
\newtheorem{lemma}{Lemma}[section]
\newtheorem{proposition}{Proposition}[section]
\newtheorem{corollary}{Corollary}[section]

\theoremstyle{explain}
\newtheorem{definition}{Definition}[section]
\newtheorem{example}{Example}[section]

\theoremstyle{aside}
\newtheorem{exercise}{\textcolor{blue}{Exercise}}[section]     % its own numbering: a separate list
\newtheorem*{remark}{Remark}
\newtheorem*{notation}{Notation}

\renewcommand{\qedsymbol}{\textcolor{accent}{$\blacksquare$}}

%% amsthm's proof environment, with tighter vertical space to match.
\makeatletter
\renewenvironment{proof}[1][\proofname]{\par
  \pushQED{\qed}%
  \normalfont \topsep=3pt plus 1pt minus 1pt\relax
  \trivlist
  \item[\hskip\labelsep\itshape #1\@addpunct{.}]\ignorespaces
}{%
  \popQED\endtrivlist\@endpefalse
}
\makeatother

%% ---- Writing aids ----------------------------------------------------
\newcommand{\defn}[1]{\textcolor{accent}{\textit{#1}}}      % term being defined
\renewcommand*{\marginfont}{\normalfont\footnotesize\color{quiet}}
\newcommand{\mnote}[1]{\marginnote{#1}}                      % aside in the margin
\newcommand{\todo}[1]{\marginnote{\color{accent}\textsc{to do}\\ #1}}

%% ---- Number systems --------------------------------------------------
\newcommand{\NN}{\mathbb{N}}
\newcommand{\ZZ}{\mathbb{Z}}
\newcommand{\QQ}{\mathbb{Q}}
\newcommand{\RR}{\mathbb{R}}
\newcommand{\CC}{\mathbb{C}}
\newcommand{\FF}{\mathbb{F}}

%% ---- Delimiters: \abs{x}  \norm{v}  \inner{u,v}  \set{x \mid x>0}
%% ---- (starred form, e.g. \norm*{\frac12 v}, grows to fit)
\DeclarePairedDelimiter{\abs}{\lvert}{\rvert}
\DeclarePairedDelimiter{\norm}{\lVert}{\rVert}
\DeclarePairedDelimiter{\inner}{\langle}{\rangle}
\DeclarePairedDelimiter{\set}{\{}{\}}

%% ---- Operators -------------------------------------------------------
\DeclareMathOperator{\im}{im}
\DeclareMathOperator{\rank}{rank}
\DeclareMathOperator{\tr}{tr}
\DeclareMathOperator{\Span}{span}
\DeclareMathOperator{\Hom}{Hom}
\DeclareMathOperator{\End}{End}
\DeclareMathOperator{\GL}{GL}
\newcommand{\id}{\mathrm{id}}
\newcommand{\T}{^{\top}}      % A\T       transpose
\newcommand{\inv}{^{-1}}      % A\inv     inverse
